\section{Introduction}
\label{sec:intro}

\plan{Science research articles do not use \texttt{\textbackslash section}
labels in the main narrative. The headings here exist so this draft has a usable
table of contents; strip them at submission. Paragraph scaffolding is kept
visible so collaborators can see the intended arc. Target: four paragraphs,
${\sim}800$ words.}

\subsection*{P1 --- The promise, and where it stalls}
\begin{planlist}
  \item RFdiffusion and its successors design enzyme backbones from scratch for
        chemistries with no natural template. Backbone generation is no longer
        the bottleneck.
  \item First-pass designs are catalytically weak --- typically orders of
        magnitude below a useful enzyme. The bottleneck moved downstream, to
        optimization.
  \item \textbf{The framing this paper adopts:} optimization is limited by what
        you can physically \emph{build and test}, not by what you can compute.
        That reframes the problem as library construction.
\end{planlist}
\todo{p1 --- de novo design succeeds at backbones, stalls at activity}

\subsection*{P2 --- Three properties, and why nothing delivers all three}
\begin{planlist}
  \item A library that can rescue a weak de novo enzyme has to be
        simultaneously \textbf{large} (the active region is far away),
        \textbf{diverse} (it must cover directions, not just distance) and
        \textbf{target-specific} (every member must still fold and hold the
        active site). Name these three explicitly here --- they are the spine
        of the whole paper and each results section closes one of them.
  \item Site-saturation mutagenesis is specific and buildable but strictly
        local; it cannot reach a variant 30 mutations away.
  \item Vanilla MPNN re-sampling is structurally informed but functionally
        blind. Low temperature returns one solution sampled repeatedly; high
        temperature returns diversity that does not fold. Section~4 measures
        exactly this trade-off.
  \item Scale alone is available --- \citet{weinstein2026manufacturing} reach
        ${\sim}10^{16}$ designs by making the generative model
        \emph{manufacturing-aware}, factorizing it to match what combinatorial
        DNA assembly can physically sample. \textbf{This is the closest prior
        art to Move 2 and must be cited as such, not as a foil.} The distinction
        to draw: they constrain the model to the chemistry from the outset and
        demonstrate on antibody repertoires, where a natural distribution
        supplies the objective; we train for catalytic geometry first and
        factorize afterwards, because for a de novo enzyme there is no natural
        distribution to imitate and the reward has to be built.
  \item Cross-reference Section~\ref{sec:related} for the full treatment.
\end{planlist}
\todo{p2 --- the three properties; SSM, MPNN resampling, and the scale-only
  approaches. Cite weinstein2026manufacturing here as nearest prior art}

\subsection*{P3 --- What we contribute}
\begin{planlist}
  \item \textbf{(i) A frontier shift, not a knob.} GRPO fine-tuning of
        ProteinMPNN against a structure-prediction reward, with a KL anchor to
        the base model that holds the output distribution open. Measured on a
        public benchmark against the full temperature ladder.
  \item \textbf{(ii) Combinatorial realization.} Fragment splitting and
        recombination convert the trained policy's output into a library that is
        exponentially larger than the number of DNA parts synthesized, while
        keeping the pass rate the policy bought.
  \item \textbf{(iii) A closed loop that actually closed.} Three rounds on a de
        novo PETase, each library designed using the measured activity of the
        last.
  \item Frame (i) and (ii) as \emph{method} and (iii) as \emph{demonstration}.
        Resist the temptation to present three co-equal applications --- the
        earlier draft did, and it diluted all three.
\end{planlist}
\todo{p3 --- three contributions}

\subsection*{P4 --- Headline numbers + roadmap}
\begin{planlist}
  \item Frontier: \textbf{[3.98]}$\times$ pass rate at matched diversity on the
        AME benchmark; $U_{15}$ \textbf{[808]} vs \textbf{[608]} for the best
        baseline temperature.
  \item Fragments: \textbf{[N]} parts $\rightarrow$ \textbf{[10\^{}X]} reachable
        variants at \textbf{[XX]}\,\% retained pass rate.
  \item PETase: \textbf{[250]}$\times$ over the starting design; best variant
        ${\sim}$\textbf{[30]} mutations out; exceeds characterized natural
        enzymes on the same probe.
  \item One-sentence roadmap of the three results sections.
\end{planlist}
\todo{p4 --- headline results and roadmap}
